\begin{frame}{islo completed 255/256 restore--run--capture calls, median 6.87 s.\ev{\tagM}}
\vspace{.1cm}
{\small
\begin{tabularx}{\textwidth}{@{}Y r r r@{}}
\toprule
Workflow & Success & p50 & p95\\
\midrule
Daytona: create & 1,024/1,024 & 0.20 s & 1.12 s\\
Tensorlake: create & 256/256 & 3.44 s & 9.00 s\\
islo: restore, run, capture & 255/256 & 6.87 s & 9.04 s\\
\bottomrule
\end{tabularx}\par}
\vspace{.2cm}
{\small
\begin{ticks}
\item Restore path: a 141 MB named snapshot; concurrency 12.
\item Creation paths: requested concurrency 8.
\item Each row names its timing endpoint; teardown is excluded.
\end{ticks}\par}
\sources{Eliaz, \href{https://arxiv.org/abs/2607.09689}{arXiv:2607.09689}, Table 2.}
\qadetail{Daytona and Tensorlake requested concurrency eight. islo used a named 141 MB snapshot at concurrency twelve. Teardown is excluded from per-operation percentiles. The islo result is a client API round trip including restore, command execution and capture; memory preservation is unverified, so the supported description is restore fan-out. These are exercised paths, not isolated mechanism latencies or a provider ranking. Operations, workloads, concurrency and timing endpoints differ. The controlled reuse experiment compares equivalent usable states against a warm baseline.}
\note{[0:55] The islo trace completed two hundred fifty-five of two hundred fifty-six restore, run and capture calls. Its median was six point eight seven seconds, and its ninety-fifth percentile was nine point zero four seconds. The input was a named hundred forty-one megabyte snapshot, with concurrency twelve. Daytona and Tensorlake exercised creation at requested concurrency eight; their endpoints are shown separately. Teardown is outside these percentiles. Memory preservation still needs the nonce probe. These traces establish concrete API paths and costs; the next experiment asks whether that reuse path beats a faithful warm alternative.}
\end{frame}
